\documentclass[../../../main.tex]{subfiles}
\begin{document}

\subsection{Forest Plots}

The most common way to visualize meta-analyses is through forest plots. See \autoref{Figure:ForestPlot} for an example created from a dataset containing studies examining the effect of so-called ``third wave'' psychotherapies on perceived stress in college students. For each study, the standardized mean difference between a treatment and control group at post-test was calculated, and Hedges' $g$ small sample correction was applied \autocite{Harrer2021}.

For each study, a graphical representation of the effect size is provided. This visualization shows the point estimate of a study on the horizontal axis. This point estimate is supplemented by a line, which represents the range of the confidence interval calculated for the observed effect size. Usually, the points estimate is surrounded by a square. The size of this square is determined by the weight of the effect size: Studies with a larger weight are given a larger square \autocite{Harrer2021}.

At the bottom of the plot, a diamond shape represents the average effect. The length of the diamond symbolizes the confidence interval of the pooled result on the horizontal axis. Typically, forest plots also include a vertical reference line, which indicates the point on the horizontal axis equal to no effect \autocite{Harrer2021}. 

\begin{figure}[!htbp]
\centering
\includegraphics[width=0.8\linewidth]{appendix/methods/analysis/assets/ma/forest.png}
\caption[Third Wave forest plot]{Example of a forest plot from the Third Wave dataset.}
\label{Figure:ForestPlot}
\end{figure}
\end{document}
